\section{S-矩阵}
定义S-矩阵：
    \begin{equation}
    S_{\beta\alpha} = \braket{\beta;- | \alpha;+},
\end{equation}
表示入态 $\ket{\alpha;+}$ 演化到出态 $\ket{\beta;-}$ 的概率幅。


\begin{center}
    \vspace{0.5cm}
        \text{\color{red}\Large ***}\\
    \vspace{0.5cm}
\end{center}
这里来推导一下,S矩阵散射振幅表示:\color{red}记得写\normalcolor
\begin{center}
    \vspace{0.5cm}
        \text{\color{red}\Large ***}\\
    \vspace{0.5cm}
\end{center}

\begin{tcolorbox}[colback=white!5,colframe=black!55,title=渐进形式]
    S-算符是：
\begin{equation}
    S=\Omega^\dagger(\infty)\Omega(-\infty)=U(-\infty,\infty)
\end{equation}
其中\begin{equation}\label{5.5.2}
   \begin{aligned}
    U(t,t_0)=&\Omega^\dagger(t)\Omega(t_0)=e^{iH_0t}e^{-iH(t-t_0)}e^{-iH_0t_0}\\
   \overset{t\to\infty,t_0\to-\infty}{\longrightarrow}
   &U(\infty,-\infty)=T \exp\!\left(-i\int_{-\infty}^{\infty} dt\,V(t)\right)
   \end{aligned}
\end{equation}
在这种情况下，$S$矩阵变为\footnote{\color{blue}值得注意的是，
    ：Dyson 级数虽是渐进展开计算(并没有限定这个相互作用的大小，但仅仅在微扰的时候才能有效) S 矩阵的标准方法，但其依赖的相互作用绘景在严格数学意义下受 \textbf{Haag 定理} 限制（即相互作用表象在无限自由度下可能不存在）。
    之后公理化 QFT 引入 \textbf{海森堡场 (Heisenberg Field)} 则是为了建立不依赖于微扰展开的严格定义。
    物理上，严格的 S 矩阵元素是通过 \textbf{LSZ 约化公式} 从海森堡场的格林函数中提取的，而 Dyson 级数实则是对这一过程的渐近近似计算。
}
\begin{equation}
    S_{\beta\alpha}\to\braket{\beta|T \exp\!\left(-i\int_{-\infty}^{\infty} dt\,V(t)\right)|\alpha}
\end{equation}
\end{tcolorbox}
